\documentclass[12pt,a4paper]{article}
\usepackage{fontspec}
\setmainfont{Times New Roman}
\usepackage[russian,english]{babel}
\usepackage{geometry}
\geometry{top=18mm,bottom=20mm,left=20mm,right=20mm}
\usepackage{graphicx}
\usepackage{booktabs}
\usepackage{array}
\usepackage{microtype}
\usepackage{caption}
\usepackage{xcolor}
\usepackage{hyperref}
\usepackage{enumitem}
\usepackage{amsmath}
\usepackage{siunitx}
\sisetup{output-decimal-marker={,},detect-all}
\graphicspath{{figures/}}
\hypersetup{colorlinks=true,linkcolor=black,urlcolor=blue,citecolor=black}
\captionsetup{font=small,labelfont=bf,labelsep=period}
\setlength{\parindent}{1.25cm}
\setlength{\parskip}{0pt}
\setlength{\tabcolsep}{5pt}
\renewcommand{\arraystretch}{1.12}
\emergencystretch=2em
\newcommand{\ci}[2]{[#1;~#2]}

\begin{document}
\selectlanguage{russian}

\begin{center}
{\bfseries\large КАЧЕСТВО В ПРОСТРАНСТВЕ ПРИЗНАКОВ НЕ ГАРАНТИРУЕТ ПОЛЕЗНОСТЬ СИНТЕТИЧЕСКИХ ДЕРМАТОСКОПИЧЕСКИХ ИЗОБРАЖЕНИЙ ПРИ ДИСБАЛАНСЕ КЛАССОВ}\par
\vspace{0.7em}
{\bfseries Н.К. Захаров, И.А. Матвеева}\par
Самарский университет, Самара, Россия\par
\texttt{nk.z4kharov@gmail.com}
\end{center}

\noindent\textbf{Аннотация.}
Синтетические изображения часто применяются для компенсации дисбаланса классов, однако их оценивают без равного по объёму контроля на повторном предъявлении реальных данных. В работе исследовано, добавляет ли синтез, отобранный по геометрии пространства признаков, полезную информацию относительно согласованного по источникам повторного предъявления реальных изображений. HAM10000 разделён по группам поражений без пересечения между обучением и валидацией. ConvNeXt-S при разрешении $384\times384$ обучался с одинаковой инициализацией и вычислительным бюджетом при трёх случайных инициализациях. Для оценки использованы макро-AUPRC, макро-F1, MCC, калибровка и иерархический бутстреп по группам поражений (5000 повторов). В разведочном анализе эффект зависел от геометрического слоя: синтетика из строгой внутридоменной области увеличила MCC на 0,066, тогда как случайный остаток уменьшил MCC на 0,049. На более сильной подтверждающей базовой модели рост макро-F1 на 0,012 сопровождался снижением макро-AUPRC на 0,014 и AUPRC для меланомы на 0,045. Диагностика показала меньшие точность, плотность и покрытие в 11 из 12 срезов «кодировщик--класс», несмотря на более высокое разнообразие в 11 из 12 срезов. Причинное разложение установило, что полный конвейер «кадрирование--VAE--один шаг удаления шума» уменьшил MCC на 0,029; значительная часть потери качества ранжирования возникала уже при предварительном кадрировании. Результаты показывают, что визуальная правдоподобность и близость в одном представлении базовой модели недостаточны для вывода о прикладной полезности.

\noindent\textbf{Ключевые слова:} дермоскопия, синтетические данные, дисбаланс классов, пространство признаков, классификация изображений.

\section*{1. Введение}

Дисбаланс диагнозов в дерматоскопических наборах данных приводит к тому, что высокая общая точность может сочетаться с низкой чувствительностью редких классов. Простые методы --- взвешивание функции потерь, повторная выборка с увеличением или уменьшением частоты классов, корректировка логитов --- меняют частоту предъявления классов, но не создают новых наблюдений. Генеративные модели потенциально расширяют поддержку редких классов, однако визуальное качество синтетики не означает, что она несёт новую дискриминативную информацию для классификатора.

В опубликованных исследованиях диффузионная аугментация нередко сравнивается с естественной выборкой, но не с равным по дозе повторным предъявлением тех же реальных источников. Это создаёт смешение двух эффектов: дополнительного числа обучающих шагов и собственно нового содержания синтетики. Кроме того, FID, расстояние в DINOv2, точность и полнота генератора или субъективная реалистичность не являются прямыми оценками качества ранжирования диагнозов [1--5]. Для медицинских данных особенно важны сохранение диагноза, отсутствие утечки между изображениями одного поражения и раздельный анализ метрик ранжирования, порогового решения и калибровки.

Цель работы --- проверить, гарантирует ли близость синтетических дерматоскопических изображений к реальным данным в пространстве признаков их добавочную полезность относительно повторного предъявления реальных изображений, согласованного по источникам и количеству. Вклад работы состоит из четырёх частей: (1) защищённого от утечки протокола по группам поражений; (2) стратифицированного по геометрии сравнения синтетики с реальным контролем; (3) диагностики несогласованности представлений базовой и прикладной моделей; (4) причинного разложения предварительной обработки, VAE и удаления шума в генеративном конвейере.

\section*{2. Материалы и методы}

\subsection*{2.1. Данные и защита от утечки}

Основные эксперименты проведены на HAM10000 [6], содержащем семь диагностических классов. Единицей разделения был идентификатор поражения \texttt{lesion\_id}: все изображения одного поражения попадали только в одну часть данных. Источники синтетических изображений также исключались из валидации только на реальных данных. Заблокированная тестовая часть не использовалась ни для выбора признакового пространства, ни для настройки генератора, порогов или эпохи; приведённые ниже результаты являются внутренней валидационной оценкой.

В качестве дополнительной проверки переносимости протокола данных использован ISIC 2019: после изоляции четырёх точных дубликатов с конфликтующими метками осталось 25,327 изображений и 13,926 связанных групп. Разбиение на обучение, валидацию и заблокированную тестовую часть составило 17,729/3,799/3,799 наблюдений. Эти результаты не объединялись с HAM10000 как единая складка перекрёстной проверки.

\subsection*{2.2. Классификатор и контролируемое сравнение}

Основная подтверждающая базовая модель --- ConvNeXt-S [7], предварительно обученная на ImageNet, с входом $384\times384$. Для каждого сравнения использовались случайные инициализации 42, 43 и 44, одинаковые начальные веса, одно разбиение данных, одинаковые число предъявлений, расписание оптимизатора и критерий ранней остановки. Более тяжёлые архитектуры ConvNeXt-B и DINOv2-B/14 [8] предварительно квалифицировались как базовые модели, обучаемые только на реальных данных; ConvNeXt-S выбран как наиболее устойчивый компромисс качества и вычислительной стоимости.

Согласованный по источникам контроль повторно предъявлял реальные изображения, от которых были получены синтетические образцы. Поэтому разность между синтетикой и реальным контролем измеряет добавочную информацию преобразования, а не дополнительную дозу данных. Повторная выборка с увеличением и уменьшением частоты классов, взвешенная перекрёстная энтропия, Balanced Softmax и cRT использовались как вычислительно дешёвые методы контроля дисбаланса [9--11].

\subsection*{2.3. Генерация и отбор}

Исторический генератор использовал латентную диффузионную модель Stable Diffusion 1.5 [12] в режиме преобразования изображения. Кандидаты для классов меланомы (mel), актинического кератоза и внутриэпителиальной карциномы (akiec), а также доброкачественных кератозоподобных поражений (bkl) отбирались по расстояниям в пространствах признаков DINOv2 и прикладной ConvNeXt. Геометрические слои задавались до оценки классификатора:
\begin{itemize}[nosep,leftmargin=1.8em]
  \item \texttt{strict\_id} --- ближайшая внутридоменная область;
  \item \texttt{aid\_radial} --- умеренное радиальное расширение поддержки;
  \item \texttt{ood\_far} --- удалённые кандидаты;
  \item \texttt{random\_remaining} --- случайный остаток пула.
\end{itemize}

Для многомодельной диагностики вычислялись улучшенные точность и полнота, плотность и покрытие (PRDC), а также мера разнообразия Vendi [2--4]. Отдельная квалификация генератора на ISIC 2019 проверяла согласованность диагноза, IoU маски, ошибку вне маски и отсутствие близких дубликатов до запуска классификатора.

\subsection*{2.4. Метрики и статистический анализ}

Основной показатель качества ранжирования --- макро-AUPRC в схеме «один класс против остальных». Дополнительно рассчитаны макро-F1, MCC, сбалансированная точность, ECE, наименьшая полнота среди классов, AUROC/AUPRC каждого класса, точность, полнота, F1 и AUPRC для меланомы, а также показатели при фиксированной специфичности. Эффекты вычислялись попарно при одинаковых случайных инициализациях. Доверительные интервалы строились 5000-кратным иерархическим бутстрепом с повторной выборкой групп поражений, а не отдельных изображений. Такой анализ отделяет изменение ранжирования от изменения рабочей точки классификатора.

\section*{3. Результаты}

\subsection*{3.1. Геометрия модерирует полезность, но не гарантирует её}

На разведочном этапе 10 знак эффекта зависел от положения синтетики в пространстве признаков (рис. 1). Для \texttt{strict\_id} MCC увеличился на 0,066, 95\% ДИ \ci{0,005}{0,125}. Для \texttt{random\_remaining} MCC снизился на 0,049, 95\% ДИ \ci{-0,086}{-0,013}. Промежуточный слой \texttt{aid\_radial} дал малый неопределённый эффект, а удалённая синтетика \texttt{ood\_far} имела отрицательную тенденцию. Следовательно, неразборчивая генерация может быть вреднее согласованного реального контроля, а близость является модератором, но не достаточным критерием.

\begin{figure}[ht]
\centering
\includegraphics[width=0.96\textwidth]{figure_1_stage10_geometry_forest.png}
\caption{Парный эффект синтетики относительно согласованного по источникам реального контроля для геометрических слоёв этапа 10. Точки --- средние по трём случайным инициализациям; интервалы --- 95\% иерархический бутстреп по группам поражений, 5000 повторов.}
\end{figure}

Положительный сигнал порогового решения \texttt{strict\_id} не перенёсся на более сильный ConvNeXt-S (табл. 1). Макро-F1 вырос на 0,0115 при всех трёх случайных инициализациях, но макро-AUPRC снизился на 0,0136 во всех трёх случаях, а AUPRC для меланомы --- на 0,0447, 95\% ДИ \ci{-0,0885}{-0,0009}. ECE ухудшился на 0,0235. Это указывает на изменение рабочей точки без улучшения качества ранжирования.

\begin{table}[ht]
\centering
\caption{Подтверждающий этап 11B: эффект \texttt{strict\_id} синтетики относительно согласованного реального контроля.}
\begin{tabular}{lrrl}
\toprule
Метрика & $\Delta$ & Победы & 95\% ДИ \\
\midrule
Макро-F1 & +0,0115 & 3/3 & [-0,0089; +0,0320] \\
MCC & +0,0096 & 2/3 & интервал включает 0 \\
Макро-AUPRC & -0,0136 & 0/3 & [-0,0336; +0,0063] \\
AUPRC меланомы & -0,0447 & 0/3 & [-0,0885; -0,0009] \\
ECE (меньше лучше) & +0,0235 & 0/3 & интервал включает 0 \\
\bottomrule
\end{tabular}
\end{table}

\subsection*{3.2. Несогласованность пространств признаков}

Этап 12 показал, что простое вырождение разнообразия не объясняет отрицательный перенос. Мера разнообразия Vendi для синтетики была выше реального контроля в 11 из 12 комбинаций «кодировщик--класс». Одновременно точность, плотность и покрытие были ниже в 11 из 12 комбинаций. В DINOv2 среднее отношение расстояния новизны к реальной окрестности составляло 0,65--0,76, то есть образцы выглядели похожими на источники. В прикладной ConvNeXt-S это отношение составляло 1,08--1,86, что указывало на выход за локальную поддержку класса.

Частотный разрыв подтвердился в 14 из 15 парных интервалов «класс--метрика». Для меланомы нижний квантиль оценок положительного класса уменьшился на 0,0236, а верхний квантиль оценок отрицательного класса вырос на 0,1127. Таким образом, синтетика одновременно ослабляла трудные положительные образцы и усиливала наиболее опасные отрицательные классы-соперники. Многомодельный алгоритм отбора позднее улучшил средние точность на 0,308 и плотность на 0,432, но макро-AUPRC классификатора снова снизился на 0,0136 относительно реального контроля. Улучшение PRDC поэтому не является достаточной целевой функцией.

\subsection*{3.3. Причинное разложение генеративного пайплайна}

Аудит показал, что 87 из 90 исторических изображений были сгенерированы при силе преобразования 0,05 и 30 шагах, фактически примерно с одним шагом удаления шума. Для локализации источника потери информации были заранее определены четыре варианта: A --- повторное предъявление исходного изображения; B --- предварительное квадратное кадрирование; C --- кадрирование и прямое-обратное преобразование VAE; D --- кадрирование, VAE и один шаг UNet.

\begin{figure}[ht]
\centering
\includegraphics[width=0.96\textwidth]{figure_2_stage15a_causal_decomposition.png}
\caption{Причинное разложение этапа 15A. Отрицательные значения означают ухудшение относительно предыдущего компонента или полного реального контроля. Все эффекты усреднены по трём парным случайным инициализациям.}
\end{figure}

Полный конвейер D--A снизил макро-AUPRC на 0,0083, макро-F1 на 0,0249 и MCC на 0,0288; для MCC 95\% ДИ составил \ci{-0,0578}{-0,0004}. Предварительное кадрирование B--A объяснило основную часть потери качества ранжирования ($\Delta$ макро-AUPRC $=-0,0064$). Прямое-обратное преобразование VAE C--B было почти нейтрально по ранжированию ($-0,0003$), а один шаг UNet D--C не добавил информации для ранжирования и ухудшил макро-F1 на 0,0193. Следовательно, обнаруженную неудачу нельзя приписать только диффузионной модели: предварительная обработка уже меняла полезную информацию.

\begin{table}[ht]
\centering
\caption{Этап 15A: средние причинные эффекты по трём случайным инициализациям.}
\begin{tabular}{lrrr}
\toprule
Сравнение & Макро-AUPRC & Макро-F1 & MCC \\
\midrule
Кадрирование $-$ контроль (B--A) & -0,0064 & -0,0084 & -0,0154 \\
VAE $-$ кадрирование (C--B) & -0,0003 & +0,0028 & -0,0030 \\
Один шаг UNet $-$ VAE (D--C) & -0,0016 & -0,0193 & -0,0104 \\
Полный конвейер $-$ контроль (D--A) & -0,0083 & -0,0249 & -0,0288 \\
\bottomrule
\end{tabular}
\end{table}

\subsection*{3.4. Переносимость протокола и квалификация генератора}

На ISIC 2019 Balanced Softmax относительно естественной выборки с перекрёстной энтропией увеличил сбалансированную точность на 0,0340, 95\% ДИ \ci{0,0153}{0,0548}, но макро-AUPRC вырос только на 0,0011 при интервале, включающем ноль; макро-F1 снизился на 0,0161, ECE ухудшился на 0,0871, AUPRC для меланомы снизился на 0,0119 при всех трёх случайных инициализациях. Это воспроизводит различие между пороговыми эффектами и качеством ранжирования уже на более крупном восьмиклассовом наборе.

Перед обучением классификатора генераторы были подвергнуты квалификации только на обучающих данных (рис. 3). Сегментатор прошёл порог допуска (Dice 0,923; IoU 0,864), однако ни исторический SD1.5, ни универсальное заполнение области, ни варианты с доменной LoRA не достигли заранее заданной согласованности диагноза 0,80. Историческое преобразование изображения показало наилучшее значение 0,625, а варианты с LoRA --- 0,313--0,333. Они не были допущены к оценке классификатора. Этот отрицательный контроль показывает необходимость измерять сохранение диагноза до дорогостоящего обучения.

\begin{figure}[ht]
\centering
\includegraphics[width=0.78\textwidth]{figure_3_generator_qualification.png}
\caption{Квалификация генераторов только на обучающих данных ISIC 2019. Пунктирные линии задают заранее установленные пороги согласованности диагноза 0,80 и IoU маски 0,70. Ни один вариант не прошёл полный набор критериев.}
\end{figure}

\section*{4. Обсуждение}

Полученные результаты не означают, что синтетические данные в дермоскопии бесполезны. Они показывают более узкое и проверяемое утверждение: для исследованных конвейеров SD1.5 преобразования/заполнения изображения и доменной LoRA визуальная правдоподобность, близость в DINOv2 и даже улучшение части многомодельных PRDC не гарантируют преимущества над согласованным повторным предъявлением реальных изображений. Такой контроль критичен, поскольку повторное предъявление редких реальных наблюдений само по себе является сильной и дешёвой стратегией.

Стратифицированный по геометрии результат этапа 10 научно содержателен, хотя величина эффекта умеренна. Положительный MCC для \texttt{strict\_id} и отрицательный для случайного остатка демонстрируют взаимодействие качества кандидатов и прикладной полезности. Однако этап 11B показывает, что отбор по одному пространству может оптимизировать пороговое решение и одновременно ухудшать ранжирование меланомы. Поэтому будущий алгоритм отбора должен обучаться на прикладном показателе полезности, учитывать неопределённость и покрытие границ классов, а не только расстояние до реального многообразия.

Причинное разложение выявило исправимую инженерную проблему: квадратное кадрирование нельзя считать нейтральным компонентом генерации. Следующие генераторы должны получать вход с сохранением соотношения сторон или с направлением по маске, контролировать неизменность области вне поражения и использовать порог сохранения диагноза. Более сильную условную по классу диффузионную модель или заполнение по маске имеет смысл запускать только после воспроизводимой квалификации на обучающих данных; увеличение длительности обучения LoRA само по себе не гарантирует сохранение метки.

\section*{5. Ограничения}

Все основные эффекты получены на внутренней валидационной части HAM10000; заблокированная тестовая часть оставалась закрытой. Три случайные инициализации позволяют оценить устойчивость знака и парную дисперсию, но не заменяют повторные групповые разбиения. Генеративные выводы относятся к конкретным конвейерам SD1.5 и LoRA и не обобщаются на все современные диффузионные модели и модели потоков. Согласованность диагноза оценивалась моделью-аудитором, а не независимой панелью дерматологов; поэтому клиническая реалистичность не заявляется. ISIC 2019 использован для переноса протокола и порога допуска генератора, а не как независимое подтверждение сравнения на HAM10000.

\section*{6. Заключение}

В контролируемой классификации дерматоскопических изображений с длиннохвостым распределением качество синтетики должно определяться её добавочной полезностью относительно согласованного повторного предъявления реальных изображений. Близость к реальным данным в представлении базовой модели и улучшение генеративных метрик не являются достаточными косвенными критериями. Геометрия влияет на знак эффекта, но положительный сигнал порогового решения может сопровождаться ухудшением AUPRC и калибровки. Причинный анализ показал, что часть потерь возникает ещё при предварительной обработке, а не только в VAE или при удалении шума. Практический протокол должен объединять групповое разделение по поражениям, согласованные контроли, метрики ранжирования, иерархический бутстреп и квалификацию сохранения диагноза до запуска классификатора.

\section*{Литература}

\footnotesize
\begin{enumerate}[label={[\arabic*]},leftmargin=2.2em,itemsep=2pt]
\item Li Y., et al. Diffusion-based data augmentation for image recognition. arXiv:2603.08364, 2026.
\item Kynkäänniemi T., Karras T., Laine S., Lehtinen J., Aila T. Improved precision and recall metric for assessing generative models. NeurIPS. 2019;32.
\item Naeem M.F., Oh S.J., Uh Y., Choi Y., Yoo J. Reliable fidelity and diversity metrics for generative models. ICML. 2020:7176--7185.
\item Friedman D., Dieng A.B. The Vendi Score: a diversity evaluation metric for machine learning. TMLR. 2023.
\item Dombrowski M., et al. The learnability gap in medical latent diffusion. arXiv:2605.17087, 2026.
\item Tschandl P., Rosendahl C., Kittler H. The HAM10000 dataset: a large collection of multi-source dermatoscopic images of common pigmented skin lesions. Sci Data. 2018;5:180161. \url{https://doi.org/10.1038/sdata.2018.161}.
\item Liu Z., Mao H., Wu C.-Y., Feichtenhofer C., Darrell T., Xie S. A ConvNet for the 2020s. Proc. CVPR. 2022:11976--11986. \url{https://doi.org/10.1109/CVPR52688.2022.01167}.
\item Oquab M., et al. DINOv2: learning robust visual features without supervision. TMLR. 2024.
\item Menon A.K., Jayasumana S., Rawat A.S., Jain H., Veit A., Kumar S. Long-tail learning via logit adjustment. ICLR. 2021.
\item Kang B., Xie S., Rohrbach M., Yan Z., Gordo A., Feng J., Kalantidis Y. Decoupling representation and classifier for long-tailed recognition. ICLR. 2020.
\item Yang Y., et al. A comparison of random oversampling and undersampling for imbalanced data classification. J Big Data. 2023;10. \url{https://doi.org/10.1186/s40537-023-00857-7}.
\item Rombach R., Blattmann A., Lorenz D., Esser P., Ommer B. High-resolution image synthesis with latent diffusion models. Proc. CVPR. 2022:10684--10695. \url{https://doi.org/10.1109/CVPR52688.2022.01042}.
\item Jiang X., et al. Synthetic data generation for long-tail medical image classification: a case study in skin lesions. arXiv:2605.03221, 2026.
\item Ma C., Lyu L., Zhang G. When does synthetic data augmentation improve score-based imbalanced classification? arXiv:2606.26053, 2026.
\end{enumerate}
\normalsize

\vfill
\noindent\footnotesize\textit{Статус рукописи: черновик для конференции, 31.07.2026. Все численные результаты получены из версионируемых артефактов JSON/CSV. Заблокированная тестовая часть не открывалась. Перед отправкой требуется подтверждение транслитерации имени, списка авторов и согласия соавтора.}

\end{document}
